% !TeX root = ../main.tex
% !TeX encoding = UTF-8

\chapter{The Ledger That Was Not There, Financial Services and the 2008 Information Failure}

The question that recurs through every serious post-mortem of the 2007 to 2009 global financial crisis is not whether the risks were present. They were. The structured credit products, the funding models, the leverage ratios, the counterparty exposures, all of it was, in retrospect, visible. Economists and risk managers who examined the data available at the time have concluded, with reasonable consistency, that the information from which the crisis could have been predicted existed in the system. What did not exist was the information picture: the coherent, contemporaneous account of what each institution's decision-makers knew, at each specific decision moment, and what they did with it.

The Financial Crisis Inquiry Commission heard testimony from the chief executives, risk officers, and board members of the institutions at the centre of the crisis. The pattern of testimony was, across institution after institution, a variant of the same admission: the witness could not now reconstruct, with contemporaneous precision, what they had known at specific decision moments. Not that they had not been informed, they had. Not that the risk did not exist, it did. But the specific information picture at the specific moment of each consequential decision was not available to them as a coherent, retrievable, contemporaneous record. The inquiry was left to ask not ``what did you know?'' but ``what should you have known?'', because the evidence to answer the first question had not been captured. The second question is a standard of care question. The first is an evidentiary one. The crisis produced the regulatory architecture to address the former. The evidentiary question, what was the information picture at the decision moment, remained unanswered.

The FCIC's 2011 report found, at institution after institution, the same structural failure: boards and senior executives who had approved products, strategies, and exposures on the basis of information pictures that were incomplete, unretained, or irrecoverable by the time the examinations began. The information existed. The information picture, contemporaneous, specific, and tied to each decision moment, did not.

This chapter examines three cases from the financial crisis era in which the Digital Alibi gap produced not merely a governance failure but an evidentiary one: a situation in which institutions and individuals could not demonstrate, from contemporaneous evidence, what they had known and when they had acted. The cases are Northern Rock, Citigroup, and Goldman Sachs. They are chosen not because they are the most egregious failures of the period, they are not, but because they illustrate, with precision, three distinct forms of the same underlying gap.

\scenebreak

\section{The Financial Services Digital Alibi Gap: Three Distinct Forms}

Before examining each case, it is worth stating what distinguishes the financial services context from the aviation context examined in Chapter 5. In aviation, the Digital Alibi gap was a genuine technological and regulatory gap: the information capture technology existed but was not mandated, and no institution had built a governance architecture designed to preserve the information picture at the decision moment. The gap was structural and shared across the industry.

In financial services, the gap took a more complex form. Some institutions had information capture systems that were, in principle, capable of preserving the decision-moment information picture. The gap was not always technological. It was frequently architectural, a governance design that captured what an institution intended to record rather than what existed at the decision moment. In some cases, the gap was the result of institutional choices about what to retain and what to allow to expire. In others, it was the result of information flowing into a decision without leaving any trace in the governance record that a decision had been made on the basis of that specific information.

These three forms of the financial services Digital Alibi gap, technological, architectural, and intentional, map directly onto the three forms most likely to be found in AI governance today. The institutions that deployed AI in the 2020s without output-moment capture are replicating, in a new technological domain, the same governance failure that financial services institutions produced in the 2000s with structured credit. The technology is different. The gap is the same.

\scenebreak

\section{Northern Rock: The Board That Could Not Show What It Knew}

\textbf{The Decision That Could Not Be Reconstructed} Northern Rock plc, 2007

Northern Rock was a British mortgage bank that had, between 2001 and 2007, transformed its business model from one funded primarily by retail deposits to one funded primarily by short-term wholesale borrowing, securitising its mortgage book and using the proceeds to fund continued lending growth. By 2007, it was structurally dependent on short-term funding markets. When those markets froze in August 2007 following the collapse of the US subprime mortgage market, the bank required emergency liquidity assistance from the Bank of England and was nationalised in February 2008.

The FSA's investigation into Northern Rock's governance produced a finding that is, for the purposes of this book, more instructive than the bank's failure itself: the FSA found that the board had approved a funding model whose risk profile was knowable from the information available, but that the information picture at the moment of each board approval could not be reconstructed with forensic precision.

The board meetings at which the funding model was reviewed and approved were documented in board minutes. Those minutes recorded the decisions taken. They did not record, with the specificity the FSA's examination required, what risk analysis had been presented, in what form, at what level of granularity, at the moment of each approval. The board minutes were evidence of what was decided. They were not evidence of the information picture at the decision moment, of what the board saw, what it was not shown, and what questions it asked or did not ask before approving the model it eventually could not sustain.

\begin{quote}
\emph{The board minutes recorded the decisions. They did not record the information picture at the decision moment. The gap between those two things was the FSA's finding. It is the Digital Alibi gap.}
\end{quote}

The FSA's 2008 report on Northern Rock concluded that the board had failed to challenge management's optimistic assumptions adequately and had not ensured that the risks of its funding model were properly stress-tested. What the report could not establish with certainty, because the contemporaneous record did not exist, was what the board had actually been shown about the model's vulnerability to a wholesale funding disruption, and when. The finding was based on what the board should have known, not on what the board did know. The former is a standard of care question. The latter is an evidentiary one. Without the information picture at the decision moment, only the former could be assessed.

\scenebreak

\section{Citigroup: Forty-Three Billion Dollars of Undisclosed Risk}

\textbf{The Decision That Could Not Be Distinguished} Citigroup Inc., 2007-2008

Citigroup's exposure to the US subprime mortgage market, through its holdings of collateralised debt obligation tranches and its retained interests in structured investment vehicles, was not hidden from its board. The Financial Crisis Inquiry Commission found that Citigroup's senior risk officers and board risk committee had received presentations on the bank's CDO exposure throughout 2006 and 2007. The information was present in the governance record. It had been presented at board level.

The problem identified by the FCIC was more specific and more damaging than simple non-disclosure. The board could not demonstrate, from the governance record, what it had understood about the specific risk characteristics of its CDO positions at the moment of each decision to retain or expand those positions. The presentations to the board described the positions at an aggregate level. They did not, and in many cases could not, given the complexity of the instruments, establish the information picture at the level of each specific decision: what the board understood about the correlation assumptions embedded in the models, what it had been told about the liquidity assumptions, what stress scenarios it had reviewed before approving the positions, and at what point it had accepted the risk that the stress scenarios did not capture.

By late 2007, Citigroup held approximately \$43 billion in CDO positions on its balance sheet, positions the FCIC found the board had not fully understood at the moment they were approved. The phrase ``not fully understood'' is the evidentiary gap stated in plain language. Not understood in what specific way? By which board members? At which meetings? Based on which presentations? The FCIC report was careful to avoid attributing individual culpability without a contemporaneous record that would sustain such attribution. The record was not sufficient to sustain it. The information picture at each decision moment had not been captured in a form that could distinguish between what the board had seen and what it had understood.

\begin{quote}
\emph{The presentations were given. The decisions were taken. But the information picture at each decision moment, what was understood, not merely what was shown, was irrecoverable. Governance without contemporaneous evidence of comprehension is governance in name only.}
\end{quote}

The distinction Citigroup illustrates is the most subtle and most relevant for AI governance. It is the difference between information being present in a governance process and the information picture being captured at the decision moment. Citigroup's board had received information about its CDO exposure. What it could not demonstrate was that the information it received was understood in a form that would have allowed an informed decision, rather than a ratification of management's existing position. For AI governance, this distinction is acute: a board that receives quarterly AI governance reports has received information. A board whose governance architecture captures the information picture at each output moment, what the AI system knew, what it was not given, who reviewed it, and what authority they exercised, has evidence of an information picture. These are different things.

\scenebreak

\section{Goldman Sachs and 1MDB: When the Accountability Chain Broke}

\textbf{The Decision That Could Not Be Owned} Goldman Sachs and 1Malaysia Development Berhad, 2012-2016

The 1Malaysia Development Berhad scandal represents the most expensive enforcement action arising from the 2010s financial crime era, and the most instructive case for the accountability chain element of the Digital Alibi. Between 2012 and 2013, Goldman Sachs underwrote three bond issuances for 1MDB, the Malaysian sovereign wealth fund, raising approximately \$6.5 billion. A substantial portion of those funds was subsequently misappropriated through a conspiracy involving senior Malaysian officials and Goldman employees. The United States Department of Justice enforcement action against Goldman Sachs, concluded in 2020, resulted in a \$2.9 billion settlement.

The governance question at the heart of the Goldman enforcement action was not primarily about the bond transactions themselves, the technical execution of the transactions was straightforward. It was about the accountability chain: who in Goldman's governance structure was accountable for the due diligence decisions made before and during each bond issuance, and what information picture existed at the moment each accountability decision was made.

The DOJ's statement of facts, contained in the deferred prosecution agreement, established that significant red flags about 1MDB's governance and the legitimacy of the transaction structure had been identified by Goldman's compliance function during the due diligence process. Those red flags had been documented. The question the enforcement action pursued was whether the accountability chain, the named individuals responsible for approving the transactions despite those red flags, was documented in a form that allowed the DOJ to establish who had accepted the risk at each decision moment, and on the basis of what information picture.

The settlement rather than prosecution outcome reflected, in part, the difficulty of establishing individual culpability through the governance record. The record showed that decisions had been taken. It did not always establish, with the precision an individual prosecution would require, which named individual had accepted accountability for which specific decision, at which moment, on the basis of which specific information about the red flags that had been identified. The accountability chain was present in principle. Its documentation at the decision moment was not always sufficient to sustain the attribution of individual responsibility.

The most instructive governance detail to emerge from the enforcement proceedings was what practitioners call shadow accountability, the gap between the formal accountability map and the operational reality of who was actually making the decisions. Goldman's formal governance structure assigned accountability for the 1MDB transactions through documented approval chains and risk committee sign-offs. The operative decisions, the specific judgements about whether to proceed despite the compliance findings, and on what terms, were made through a parallel channel of informal communication and relationship-based decision-making that left no contemporaneous record in the formal governance architecture. The formal accountability map showed clean governance. The information picture at the moment of each operative decision was not captured in the formal record at all.

Shadow accountability is not unique to Goldman Sachs or to the 1MDB transactions. It is the normal operating reality of most large organisations: formal governance processes run alongside informal decision channels, and the operative decision is frequently taken in the informal channel before being ratified through the formal one. The Digital Alibi gap in this context is not merely that the output-moment record was not captured. It is that the formal record, the one that would be produced in an investigation, described a governance process that was not the process through which the operative decision was made. The governance record was contemporaneous with the formal ratification. It was not contemporaneous with the decision. For AI governance, this shadow accountability risk is acute wherever AI systems are deployed and governed through formal governance processes while operational decisions about which systems to deploy, how to configure them, and how to respond to their outputs are made through informal channels. The formal record shows policy compliance. The information picture at the operative decision moment does not exist in that record at all.

\begin{longtable}[]{@{}
  >{\raggedright\arraybackslash}p{(\columnwidth - 0\tabcolsep) * \real{1.0000}}@{}}
\toprule\noalign{}
\begin{minipage}[b]{\linewidth}\raggedright
\textbf{THE 1MDB ACCOUNTABILITY CHAIN FAILURE, THE DIGITAL ALIBI PARALLEL}
\end{minipage} \\
\midrule\noalign{}
\endhead
\bottomrule\noalign{}
\endlastfoot
Red flags documented: the information existed in the compliance record. \\
Decisions taken: the transactions proceeded. Someone approved them. \\
Shadow accountability gap: the operative decisions were made through informal channels that left no contemporaneous record in the formal governance architecture. The formal map showed clean accountability. The information picture at the moment of each operative decision was not in the formal record at all. \\
Formal accountability gap: where the governance record did capture decision-making, it did not establish, with decision-moment precision, which named individual accepted accountability for which specific risk, on the basis of which specific information picture. \\
Consequence: \$2.9 billion settlement. Individual prosecutions limited by the evidentiary gap. \\
The Digital Alibi lesson: accountability in principle is not accountability at the decision moment. Shadow accountability, where the operative decision is made outside the formal record, is invisible to any governance architecture that does not capture the information picture at the actual moment the decision was made, not the moment it was ratified. \\
\end{longtable}

\scenebreak

\section{The Three-Part Forensic Anatomy: Financial Services}

Across all three cases, the Chapter 5 forensic structure, what existed, what was accessible, what was missing, applies with consistent precision.

What existed: in all three cases, substantial information was present in the governance record. Northern Rock's board had received funding model presentations. Citigroup's board had received CDO exposure reports. Goldman's compliance function had documented due diligence red flags. The information was not hidden. It existed.

What was accessible: the information that existed was accessible to investigators, regulators, and enforcement bodies, but in aggregate, reconstructed, or decontextualised form. Northern Rock's board minutes could be read. Citigroup's presentation slides could be retrieved. Goldman's compliance documentation could be produced. What was accessible was the trace of the information that had flowed through the governance process.

What was missing: the information picture at each specific decision moment. Northern Rock's board minutes did not establish what the board was shown about stress scenarios at the moment it approved the funding model for the final time before the crisis. Citigroup's board records did not establish what each director understood about correlation assumptions at the moment each CDO position was approved. Goldman's governance records did not establish, with decision-moment precision, who had accepted accountability for proceeding despite the compliance findings, and what information they had at that moment.

\begin{quote}
\emph{In each case, the information existed. The trace of its passage through governance existed. What did not exist was the contemporaneous record of the information picture, who knew what, in what specific form, at the exact moment each consequential decision was made.}
\end{quote}

For AI governance, this three-part anatomy has a direct application. An organisation that deploys AI in material decision processes will have information about those decisions, logs, reports, model outputs, governance meeting records. That information will be accessible in the event of an investigation. What will be missing, unless a Digital Alibi architecture has been built, is the information picture at each output moment: what the AI system knew, what it was not given, who reviewed its output with what authority, and who was named and accountable at that precise instant. The financial services cases demonstrate that the gap between accessible information and a contemporaneous information picture is not closed by the volume of documentation. It is closed only by output-moment capture.

\scenebreak

\section{The Liability Consequence: Institutional Settlement and Personal Exposure}

The post-2008 regulatory response to the financial crisis information failures produced the governance architecture that financial services now operates under. The FCA's Senior Managers and Certification Regime, introduced in the UK from 2016, was directly shaped by the lesson that institutional accountability without named individual accountability at the decision moment had allowed the crisis to produce institutional consequences without individual ones. The SM\&CR's requirement for named personal accountability, documented contemporaneously, is the financial services equivalent of the 1958 aviation mandate: a regulatory response to a demonstrated governance gap that made the question ``who was accountable at the moment of this decision?'' unanswerable from the existing record.

The pattern across Northern Rock, Citigroup, and Goldman Sachs is the same pattern Chapter 3 identified in the legal architecture: institutional settlements are reached because individual prosecutions cannot be sustained without a contemporaneous information picture that establishes, beyond the standard required for criminal or significant civil liability, what a named individual knew and decided at a specific moment. The absence of that record protects individuals from the most severe consequences while exposing institutions to settlement costs that reflect the inability to demonstrate governance adequacy at the decision-moment level.

For AI governance, the implication is direct. When the first major AI enforcement actions involve material AI-assisted decisions, in credit, in underwriting, in hiring, in procurement, the pattern will be the same. Institutions that cannot demonstrate the information picture at each output moment will face institutional exposure. The individuals on their accountability maps will face the question of whether their named accountability was documented contemporaneously or assigned retrospectively. The SM\&CR, which Chapter 3 established as the named individual liability mechanism for financial services, was built precisely to prevent the individual accountability gap that the 2008 era produced. The AI governance gap is recreating the conditions that made the SM\&CR necessary.

\begin{quote}
\emph{The SM\&CR was built because the 2008 crisis demonstrated that institutional accountability without named individual accountability at the decision moment allowed catastrophic governance failures to produce institutional consequences without individual ones. AI governance is recreating those conditions.}
\end{quote}

\scenebreak

\section{The Architecture Financial Services Built}

The regulatory architecture that emerged from the 2007 to 2009 financial crisis represents the most comprehensive attempt in any sector to mandate the contemporaneous capture of decision-moment information for consequential decisions. The SM\&CR's named individual accountability requirement, combined with the PRA's model risk management guidance (SS1/23), DORA's ICT risk management obligations, and the Basel Committee's governance principles for banks, collectively create a governance standard that is closer to the Digital Alibi requirement than any prior financial services framework.

Closer, but not at. As Chapter 4 established, the SM\&CR creates the accountability chain framework without mandating output-moment capture of the information picture. SS1/23 creates model risk management requirements without mandating the contemporaneous record of what each model knew and was not given at each specific output moment. The financial services regulatory architecture, built in response to the 2008 crisis, has closed the gap further than any other sector. The Digital Alibi gap remains.

The lesson of financial services is not that the regulatory response was inadequate. It is that regulatory responses to information failures are, by their nature, responses to the specific form of the failure that was identified, and they do not anticipate the next form. The 2008 response addressed the accountability gap at institutional level and, through SM\&CR, at individual level for named Senior Managers. It did not address the output-moment information picture gap for AI-assisted decisions, because AI-assisted decisions at the scale and speed now emerging were not the form of consequential decision-making the 2008 crisis made visible. The gap the next financial services crisis will expose is the AI decision-moment gap. The architecture to close it is available now, before the crisis that will mandate it.

\scenebreak

\section{From Finance to Prosecution: The Next Reckoning}

Chapters 5 and 6 have examined two industries in which the Digital Alibi gap produced consequences primarily in the domain of regulatory enforcement and institutional accountability. Aviation produced a mandate. Financial services produced a regulatory framework and a personal liability regime. Neither produced the most extreme consequence the Digital Alibi gap can generate: the wrongful prosecution of innocent individuals on the basis of a system whose governance architecture prevented the information picture from being reconstructed.

Chapter 7 examines the Post Office Horizon scandal, the case in which the absence of an information picture at the moment of each governance decision to prosecute sub-postmasters produced not a regulatory settlement or an institutional fine, but the imprisonment and financial destruction of hundreds of people whose only failing was that the organisation prosecuting them could not be shown to have been wrong. The Digital Alibi gap, in its most consequential form, is not a governance problem. It is a justice problem. Chapter 7 is the case that makes that consequence unavoidable.
